
A simple operation --- randomly shuffling neurons before each training step --- outperforms a mathematically guaranteed equivariant architecture when only 100 trained models are available. With 250 models, the equivariant architecture wins decisively; flat-MLP does not reach 90\% of its own peak performance within any sampled training size up to N=1,000. This observation raises a question that the existing weight-space learning literature has not addressed: at what data scale does structural permutation equivariance produce a sample efficiency advantage over plain encoders, and does it hold uniformly across zoo sizes?

The setting is \textit{weight-space property prediction}: given a collection of trained neural networks (a model zoo), an encoder is learned to predict model properties --- such as test accuracy --- directly from network weights. A fundamental symmetry exists in this setting: permuting neurons in any hidden layer produces a functionally identical network with a different weight vector. Encoders that ignore this symmetry must learn the invariance from data; those that enforce it by construction --- equivariant encoders --- are expected to require fewer training examples to achieve the same prediction quality, particularly when labeled model zoos are small.

This expectation has motivated a line of equivariant encoder architectures: DWSNets \citep{Navon2023Equivariant}, GNN-NFN \citep{Kofinas2024Graph}, NFN \citep{Zhou2023Permutation}, Universal NFN \citep{Zhou2024Universal}, and Monomial-NFN \citep{Tran2024Monomial}. These encoders achieve state-of-the-art property prediction on model zoos. However, a critical question has remained unaddressed: \textit{by how much are they more sample-efficient than plain encoders, and does this advantage hold across all zoo sizes?}

A methodological obstacle prevents a direct answer. Prior equivariant encoder papers use private or custom train/test splits; prior plain encoder papers use different datasets. No controlled, shared-split comparison of equivariant versus plain weight-space encoders at systematically varied training set sizes has been conducted.

A theoretical result sharpens the question. Dayan, Eitan, and Maron [2026] proved that all permutation-equivariant weight-space networks are equivalent in \textit{expressivity} given sufficient data --- implying equivariant and plain encoders share the same function class ceiling. The difference must lie in \textit{sample complexity}. The magnitude of this difference, and whether it is uniform across data regimes, is what prior work has not measured.

This paper fills that gap with a controlled learning curve study on the ModelZooDataset CIFAR-10 CNN zoo [Schürholt et al., 2022], comparing three conditions at training sizes N $\in$ {100, 250, 500, 1000, full}: (1) flat-MLP, (2) flat-MLP with permutation augmentation (PermAug), and (3) GNN-NFN.

The key contributions are:

\begin{enumerate}
\item \textbf{A large sample efficiency advantage.} GNN-NFN reaches 90\% of its peak $R^2$ at $N\approx 147$ training models. Flat-MLP reaches only 84\% of its own peak at N=1,000 and requires the full dataset (~7,000 models) to reach 90\% of its peak $R^{2}=0.886$, yielding an efficiency ratio of approximately $47\times$. Even under the conservative bound that treats N=1,000 as the flat-MLP reference, the ratio is at least 6.$8\times$, well above the $2\times$ threshold required by the original hypothesis.
\end{enumerate}

\begin{enumerate}
\item \textbf{A novel data-regime crossover.} At N=100, PermAug ($R^{2}=0.138)$ outperforms structural equivariance (GNN-NFN $R^{2}=-0.016)$. At N=250, the ordering reverses: GNN-NFN (0.767) > PermAug (0.532) > flat-MLP (0.449). This crossover reveals a minimum-data threshold of approximately 150--250 models for equivariant graph encoders on the CIFAR-10 CNN zoo. This finding was not predicted by the original hypothesis and is not present in any prior weight-space learning study. It is based on a single random seed and requires multi-seed replication.
\end{enumerate}

\begin{enumerate}
\item \textbf{Mechanistic grounding.} Permutation equivariance is verified to floating-point precision for both GNN-NFN (max\_diff=1.$80\times 10^{-6}$ across 10,000 checks) and DWSNets (max\_diff=7.$45\times 10^{-9}$ across 1,000 checks), establishing the structural basis for the efficiency advantage.
\end{enumerate}

\begin{enumerate}
\item \textbf{A reusable evaluation protocol.} Shared-split learning curves at {100, 250, 500, 1000, full}, 90\%-peak efficiency ratio, bootstrap confidence intervals, and equivariance verification constitute a complete protocol applicable to future encoder comparisons.
\end{enumerate}

The remainder of this paper is organized as follows. Section 2 situates the work in three streams of prior research. Section 3 describes the methodology. Section 4 details the experimental setup. Section 5 presents results. Section 6 discusses implications and limitations. Section 7 concludes.

